% SEGMENT OLP-0502-S01
% part: intuitionistic-logic
% chapter: semantics
% section: topological-semantics

\documentclass[../../../include/open-logic-section]{subfiles}

\begin{document}

\olfileid{int}{sem}{top}

\olsection{இடவியல் பொருண்மையியல்}

இடவியல் என்ற கணிதக் கருத்துருவைப் பயன்படுத்துவது, உள்ளுணர்வுவாதத்
தருக்கத்திற்குப் பொருண்மையியல் தரும் மற்றொரு வழியாகும்.

\begin{defn}
  $X$ ஒரு கணமாக இருக்கட்டும். கீழுள்ள பண்புகளை நிறைவுசெய்யும் கணம்
  $\Top{O} \subseteq \Pow{X}$ ஆனது \emph{$X$ மீதான இடவியல்} ஆகும்.
  $\Top{O}$ இன் !!{element}s அவ்விடவியலின் \emph{திறந்த கணங்கள்}
  எனப்படுகின்றன. $X$ என்ற கணமும் $\Top{O}$ உம் சேர்ந்தது ஓர்
  \emph{இடவியல் வெளி} எனப்படுகிறது.
  \begin{enumerate}
  \item வெற்றுக் கணமும் முழு வெளியும் திறந்தவை: $\emptyset$, $X \in
    \Top{O}$.
  \item திறந்த கணங்கள் முடிவுறு வெட்டுகளின் கீழ் மூடப்பட்டுள்ளன: $U$, $V \in
    \Top{O}$ எனில், $U \cap V \in \Top{O}$.
  \item திறந்த கணங்கள் தன்னிச்சையான ஒன்றியங்களின் கீழ் மூடப்பட்டுள்ளன:
    எல்லா $i \in I$ க்கும் $U_i \in \Top{O}$ எனில், $\bigcup
    \Setabs{U_i}{i \in I} \in \Top{O}$.
  \end{enumerate}
\end{defn}

திறந்த கணங்களின் தொகுப்பைச் சூழலிலிருந்து அறிய முடியுமெனில், இடவியல்
வெளியை~$X$ என்று மட்டும் எழுதலாம்; இருப்பினும், திறந்த கணங்கள்~$X$ க்கு
அளிக்கப்பட்ட பின்னரே அது இடவியல் வெளி எனப்படலாம் என்பதைக் கவனிக்கவும்.
% SOURCE CORRECTION TA-IL-005: X with its open sets is a topological space,
% while the topology itself is the collection O.

\begin{defn}
  உள்ளுணர்வுவாதக் கூற்றுக்கணிப்புத் தருக்கத்தின் ஓர் \emph{இடவியல் மாதிரி}
  என்பது $\mModel{X} = \tuple{X, \Top{O}, V}$ என்ற மும்மை; இதில்
  $\Top{O}$ ஆனது~$X$ மீதான ஓர் இடவியல், மேலும் $V$ என்பது ஒவ்வொரு கூற்று
  மாறிக்கும்~$\Top{O}$ இலுள்ள ஒரு திறந்த கணத்தை மதிப்பளிக்கும் சார்பு.

  இடவியல் மாதிரி~$\mModel{X}$ கொடுக்கப்பட்டால், $\Prop{X}{!A}$ ஐப்
  பின்வருமாறு தொகுமுறையில் வரையறுக்கலாம்:
  \begin{enumerate}
  \item $\Prop{X}{\lfalse} = \emptyset$
  \item $\Prop{X}{p} = V(p)$
  \item $\Prop{X}{!A \land !B} = \Prop{X}{!A} \cap \Prop{X}{!B}$
  \item $\Prop{X}{!A \lor !B} = \Prop{X}{!A} \cup \Prop{X}{!B}$
  \item $\Prop{X}{!A \lif !B} = \Interior{(X \setminus \Prop{X}{!A}) \cup
    \Prop{X}{!B}}$
  \end{enumerate}
  இங்கு, $\Interior{V}$ என்பது $V \subseteq X$ என்ற கணத்தை அதன்
  \emph{உட்பகுதிக்கு}, அதாவது அது கொண்டுள்ள எல்லா திறந்த கணங்களின்
  ஒன்றியத்திற்கு அனுப்பும் சார்பு. வேறு சொற்களில்,
  \[
  \Interior{V} = \bigcup \Setabs{U}{U \subseteq V \text{ மற்றும் } U \in \Top{O}}.
  \]
\end{defn}

எந்தக் கணத்தின் உட்பகுதியும் திறந்த கணங்களின் ஒன்றியம் என்பதால், அது எப்போதும்
திறந்தது என்பதைக் கவனிக்கவும். எனவே $\Prop{X}{!A}$ எப்போதும் ஒரு திறந்த கணம்.

இடவியல் பொருண்மையியல் மிகவும் அருவமானதாக இருந்தாலும், அதற்கான உந்துதலைத் தரும்
வகையில் அதைப் புரிந்துகொள்ள சில வழிகள் உள்ளன. $X$ இன் !!{element}s, அல்லது
``புள்ளிகள்'', கூற்றுகளை மதிப்பிடக்கூடிய புள்ளிகள் என்று கொள்க. $!A$ மெய்யாகும்
எல்லாப் புள்ளிகளின் கணமே~$!A$ வெளிப்படுத்தும் கூற்று. புள்ளிகளின் ஒவ்வொரு
கணமும் சாத்தியமான கூற்று அல்ல; $\Top{O}$ இன் !!{element}s மட்டுமே அவ்வாறு
இருக்கும். எல்லா~$X$ க்கும் $!A$ மெய்யாகும் ஒவ்வொரு புள்ளியிலும்~$!B$
மெய்யாகும்போது, அதாவது $\Prop{X}{!A} \subseteq \Prop{X}{!B}$ ஆகும்போது,
அப்போதும் அப்போது மட்டுமே, $!A \Entails !B$. பொருத்தமற்ற கூற்று~$\lfalse$
ஒருபோதும் மெய்யன்று; எனவே $\Prop{X}{\lfalse} = \emptyset$.

உள்ளுணர்வுவாத முறையில் செல்லுபடியாகும் விதிகள் மெய்யாக இருக்க, அதாவது
$!A \Proves !B$ எனில், எனில் மட்டுமே, $\Prop{X}{!A} \subseteq
\Prop{X}{!B}$ ஆக இருக்க, $!B \land !C$, $!B \lor !C$, $!B \lif !C$
ஆகியவை வெளிப்படுத்தும் கூற்றுகள்~$!B$,~$!C$ ஆகியவை வெளிப்படுத்தும்
கூற்றுகளுடன் எவ்வாறு தொடர்புடையவையாக இருக்க வேண்டும்? எந்த~$!A$ க்கும்
$\lfalse \Proves !A$ தேவை; எல்லா $U$ க்கும் $\emptyset \subseteq U$
என்பதால் இது நிறைவுசெய்யப்படுகிறது. $!B \land !C \Proves !B$ என்பதால்,
$\Prop{X}{!B \land !C} \subseteq \Prop{X}{!B}$ தேவை; இதேபோல்
$\Prop{X}{!B \land !C} \subseteq \Prop{X}{!C}$ தேவை. $W \subseteq U$,
$W \subseteq V$ ஆகியவற்றை நிறைவுசெய்யும் மிகப்பெரிய கணம் $U \cap V$.
மறுபுறம், $!B \Proves !B \lor !C$ மற்றும் $!C \Proves !B \lor !C$;
எனவே $\Prop{X}{!B} \subseteq \Prop{X}{!B \lor !C}$ மற்றும்
$\Prop{X}{!C} \subseteq \Prop{X}{!B \lor !C}$ தேவை. $U \subseteq W$,
$V \subseteq W$ ஆகியவற்றை நிறைவுசெய்யும் மிகச்சிறிய கணம்~$W$ என்பது
$U \cup V$.
% SOURCE CORRECTION TA-IL-006: derivability corresponds to non-strict inclusion.

$\lif$ க்கான வரையறை நுட்பமானது: $!A$ உடன் சேரும்போது~$!B$ ஐ
உட்கிடையாக்கும் மிக வலுக்குறைந்த கூற்றை $!A \lif !B$ வெளிப்படுத்துகிறது.
$!A \lif !B$ ஆனது~$!A$ உடன் சேர்ந்து~$!B$ ஐ உட்கிடையாக்குகிறது என்பது
$(!A \lif !B) \land !A \Proves !B $ இலிருந்து தெளிவாகிறது. எனவே
$\Prop{X}{!A \lif !B} \cap \Prop{X}{!A} \subseteq \Prop{X}{!B}$ ஆகும்
வகையிலான மிகப்பெரிய திறந்த கணமாக $\Prop{X}{!A \lif !B}$ இருக்க வேண்டும்;
இதுவே நமது வரையறைக்கு இட்டுச்செல்கிறது.

\end{document}
